\chapter*{Notación}
\addcontentsline{toc}{chapter}{Notación}
\markboth{NOTACIÓN}{}

A lo largo de este libro utilizaremos la siguiente notación matemática.
Dado que el libro abarca múltiples áreas (análisis funcional, álgebra,
topología, geometría diferencial y aprendizaje profundo), algunos símbolos
admiten varios significados según el contexto (por ejemplo, $G$ puede denotar
un grupo o un grafo, $K$ un núcleo convolucional o una matriz de claves).
En general, el contexto matemático circundante determina la
interpretación correcta de cada símbolo.

\section*{Fundamentos Matemáticos}

\subsection*{Convenciones Notacionales}
\begin{longtable}{>{\centering\arraybackslash}p{4cm}p{10cm}}
% Primera página
\toprule
\textbf{Símbolo} & \textbf{Descripción} \\
\midrule
\endfirsthead
% Páginas siguientes
\toprule
\textbf{Símbolo} & \textbf{Descripción} {\small(continuación)} \\
\midrule
\endhead
% Pie en páginas intermedias
\midrule
\multicolumn{2}{r}{\small\textit{Continúa en la siguiente página}} \\
\endfoot
% Pie en la última página
\bottomrule
\endlastfoot
$\mathbf{x}, \mathbf{y}, \mathbf{z}$ & Vectores (negrita minúscula) \\
$\mathbf{A}, \mathbf{B}, \mathbf{C}$ & Matrices (negrita mayúscula) \\
$\mathcal{M}, \mathcal{N}$ & Variedades (caligráficas mayúsculas) \\
$G, H$ & Grupos (mayúsculas); para grafos preferir $\mathcal{G} = (V, E)$ \\
$f, g, h$ & Funciones (minúsculas) \\
$\theta, \phi, \psi$ & Parámetros o coordenadas (griegas minúsculas) \\
$\mathbb{R}$ & Números reales \\
$\mathbb{C}$ & Números complejos \\
$\mathbb{N}$ & Números naturales \\
$\mathbb{Z}$ & Números enteros \\
$x_i, w_{ij}$ & Componentes de vectores y matrices \\
$f^{(\ell)}$ & Función o capa $\ell$ en una red profunda (se usa $\ell$ para evitar confusión visual con $1$) \\
$\theta_t$ & Parámetros en el paso temporal $t$ \\
$\mathcal{M}_p$ & Estructura local en el punto $p$ \\
$h_v^{(\ell)}$ & Características del nodo $v$ en la capa $\ell$ \\
$h^{(\ell)}$ & Representación oculta en la capa $\ell$ (abreviación de $h_v^{(\ell)}$ cuando el nodo es claro del contexto) \\
$d_i$ & Grado del vértice $i$ \\
$\lambda_i$ & $i$-ésimo valor propio \\
$A \succeq 0$ & Matriz simétrica semidefinida positiva \\
$A \succ 0$ & Matriz simétrica definida positiva \\
\end{longtable}

\subsection*{Análisis Funcional y Espacios}
\begin{longtable}{>{\centering\arraybackslash}p{4cm}p{10cm}}
% Primera página
\toprule
\textbf{Símbolo} & \textbf{Descripción} \\
\midrule
\endfirsthead
% Páginas siguientes
\toprule
\textbf{Símbolo} & \textbf{Descripción} {\small(continuación)} \\
\midrule
\endhead
% Pie en páginas intermedias
\midrule
\multicolumn{2}{r}{\small\textit{Continúa en la siguiente página}} \\
\endfoot
% Pie en la última página
\bottomrule
\endlastfoot
$L^p(\Omega)$ & Espacios de Lebesgue \\
$L^2(\mathbb{R}^d)$ & Espacio de funciones de cuadrado integrable \\
$\|f\|_{L^2}$ & Norma $L^2$ de la función $f$ \\
$C^k(\mathcal{M})$ & Funciones $k$-veces diferenciables en $\mathcal{M}$ \\
$C(\mathcal{M})$ & Funciones continuas en la variedad $\mathcal{M}$ \\
$\mu$ & Medida \\
$(X, \mathcal{F}, \mu)$ & Espacio de medida (conjunto, $\sigma$-álgebra, medida) \\
$\mathcal{B}(X)$ & $\sigma$-álgebra de Borel de $X$ \\
$\lambda$ & Medida de Lebesgue \\
$\int_\Omega f \, d\mu$ & Integral de Lebesgue respecto a la medida $\mu$ \\
a.e. (c.p.d.) & Casi por doquier (almost everywhere) \\
$|\mu|$ & Variación total de la medida $\mu$ \\
$C(X)$ & Espacio de funciones continuas sobre compacto $X$ \\
$\|f\|_\infty$ & Norma supremo (uniforme) de la función $f$ \\
$X'$ & Espacio dual de $X$ (funcionales lineales continuos) \\
$\text{supp}(f)$ & Soporte de la función $f$ \\
$\text{span}(\{v_1, \ldots, v_n\})$ & Espacio vectorial generado por $\{v_1, \ldots, v_n\}$ \\
$\langle \cdot, \cdot \rangle$ & Producto interior \\
$\| \cdot \|$ & Norma euclidiana \\
$\nabla$ & Operador gradiente \\
$\nabla^2$ & Laplaciano euclidiano (operador de Laplace en $\mathbb{R}^m$) \\
$\Delta$ & Operador de Laplace-Beltrami (laplaciano continuo) \\
$L_G$ & Laplaciano discreto de un grafo $G$ \\
$f'(x)$ & Derivada de la función $f$ (notación prima) \\
$\partial$ & Derivada parcial \\
$\circ$ & Composición de funciones \\
$\arg\min$, $\arg\max$ & Argumento del mínimo/máximo \\
\end{longtable}

\subsection*{Distribuciones de Probabilidad}
\begin{longtable}{>{\centering\arraybackslash}p{4cm}p{10cm}}
% Primera página
\toprule
\textbf{Símbolo} & \textbf{Descripción} \\
\midrule
\endfirsthead
% Páginas siguientes
\toprule
\textbf{Símbolo} & \textbf{Descripción} {\small(continuación)} \\
\midrule
\endhead
% Pie en páginas intermedias
\midrule
\multicolumn{2}{r}{\small\textit{Continúa en la siguiente página}} \\
\endfoot
% Pie en la última página
\bottomrule
\endlastfoot
$\mathcal{N}(\mu, \sigma^2)$ & Distribución normal (gaussiana) con media $\mu$ y varianza $\sigma^2$ \\
$\text{Bernoulli}(p)$ & Distribución de Bernoulli con probabilidad de éxito $p$ \\
$\mathbb{E}[X]$ & Esperanza matemática de la variable aleatoria $X$ \\
$\text{Var}(X)$ & Varianza de la variable aleatoria $X$ \\
$p(x)$, $P(X=x)$ & Función de densidad/masa de probabilidad \\
$X \sim \mathcal{D}$ & La variable $X$ sigue la distribución $\mathcal{D}$ \\
$\sigma$ & Desviación estándar o parámetro de escala \\
\end{longtable}

\section*{Estructuras Algebraicas}

\subsection*{Teoría de Grupos}
\begin{longtable}{>{\centering\arraybackslash}p{4cm}p{10cm}}
% Primera página
\toprule
\textbf{Símbolo} & \textbf{Descripción} \\
\midrule
\endfirsthead
% Páginas siguientes
\toprule
\textbf{Símbolo} & \textbf{Descripción} {\small(continuación)} \\
\midrule
\endhead
% Pie en páginas intermedias
\midrule
\multicolumn{2}{r}{\small\textit{Continúa en la siguiente página}} \\
\endfoot
% Pie en la última página
\bottomrule
\endlastfoot
$H \leq G$ & $H$ es subgrupo de $G$ \\
$H \triangleleft G$ & $H$ es subgrupo normal de $G$ \\
$|G|$ & Orden (cardinalidad) del grupo $G$ \\
$[G:H]$ & Índice de $H$ en $G$ (número de clases laterales) \\
$\text{Stab}(x)$ & Subgrupo estabilizador del elemento $x$ \\
$\text{Orb}(x)$ & Órbita del elemento $x$ bajo la acción de grupo \\
$\ker(\phi)$ & Núcleo del homomorfismo $\phi$ \\
$\text{Im}(\phi)$ & Imagen del homomorfismo $\phi$ \\
$G \cong H$ & Los grupos $G$ y $H$ son isomorfos \\
$g \cdot x$ & Acción del elemento $g$ sobre $x$ \\
$C_4$ & Grupo cíclico de orden 4 \\
$D_4$ & Grupo diédrico del cuadrado (orden 8 = 4 rotaciones + 4 reflexiones) \\
$\mathbb{Z}_n$ & Grupo cíclico de enteros módulo $n$ \\
$S_n$ & Grupo de permutaciones de $n$ elementos \\
$\operatorname{sgn}(\sigma)$ & Signo (paridad) de la permutación $\sigma \in S_n$ \\
$\operatorname{Alt}$ & Operador de alternación (antisimetrización) de tensores \\
$p4, p4m$ & Grupos de simetría planar \\
\end{longtable}

\subsection*{Grupos de Lie y Álgebras de Lie}
\begin{longtable}{>{\centering\arraybackslash}p{4cm}p{10cm}}
% Primera página
\toprule
\textbf{Símbolo} & \textbf{Descripción} \\
\midrule
\endfirsthead
% Páginas siguientes
\toprule
\textbf{Símbolo} & \textbf{Descripción} {\small(continuación)} \\
\midrule
\endhead
% Pie en páginas intermedias
\midrule
\multicolumn{2}{r}{\small\textit{Continúa en la siguiente página}} \\
\endfoot
% Pie en la última página
\bottomrule
\endlastfoot
$SO(n)$ & Grupo de rotaciones en $n$ dimensiones \\
$SO(2)$ & Grupo de rotaciones en 2D (isomorfo a $U(1)$) \\
$SO(3)$ & Grupo de rotaciones en 3D \\
$O(n)$ & Grupo ortogonal en $n$ dimensiones \\
$E(n)$ & Grupo euclidiano en $n$ dimensiones \\
$SE(2), SE(3)$ & Grupo euclidiano especial en 2D y 3D \\
$GL(n, \mathbb{R})$ & Grupo de matrices invertibles $n \times n$ \\
$GL(\mathcal{F})$ & Grupo lineal general sobre espacio de características $\mathcal{F}$ \\
$SU(2)$ & Grupo especial unitario $2 \times 2$ \\
$\mathbb{H}$ & Álgebra de cuaterniones de Hamilton \\
$G/H$ & Espacio homogéneo (clases laterales de $H$ en $G$) \\
$St(k,n)$ & Variedad de Stiefel (marcos ortonormales de $k$ vectores en $\mathbb{R}^n$) \\
$Gr(k,n)$ & Grassmannian de $k$-planos en $\mathbb{R}^n$ \\
$\mathfrak{g}$ & Álgebra de Lie del grupo de Lie $G$ \\
$T_e G$ & Espacio tangente al grupo $G$ en la identidad \\
$[\cdot, \cdot]$ & Corchete de Lie \\
$\exp: \mathfrak{g} \rightarrow G$ & Mapa exponencial \\
$\mathfrak{so}(n)$ & Álgebra de Lie de $SO(n)$ (matrices antisimétricas) \\
$\mathfrak{se}(n)$ & Álgebra de Lie de $SE(n)$ \\
$\mathfrak{gl}(n, \mathbb{R})$ & Álgebra de Lie de $GL(n, \mathbb{R})$ \\
$[\mathbf{v}]_\times$ & Matriz antisimétrica asociada al vector $\mathbf{v} \in \mathbb{R}^3$ \\
$\operatorname{Ad}_g: \mathfrak{g} \to \mathfrak{g}$ & Representación adjunta del grupo \\
$\rtimes$ & Producto semidirecto de grupos \\
$\oplus$ & Suma directa de espacios vectoriales o álgebras \\
$\mu: G \times G \to G$ & Multiplicación del grupo de Lie \\
\end{longtable}

\subsection*{Representaciones de Grupos}
\begin{longtable}{>{\centering\arraybackslash}p{4cm}p{10cm}}
% Primera página
\toprule
\textbf{Símbolo} & \textbf{Descripción} \\
\midrule
\endfirsthead
% Páginas siguientes
\toprule
\textbf{Símbolo} & \textbf{Descripción} {\small(continuación)} \\
\midrule
\endhead
% Pie en páginas intermedias
\midrule
\multicolumn{2}{r}{\small\textit{Continúa en la siguiente página}} \\
\endfoot
% Pie en la última página
\bottomrule
\endlastfoot
$\rho: G \rightarrow GL(V)$ & Representación de grupo $G$ en espacio $V$ \\
$\widehat{G}$ & Conjunto de representaciones irreducibles de $G$ \\
$\mathcal{H}_\lambda$ & Espacio de representación irreducible $\lambda$ \\
$n_\lambda$ & Multiplicidad de la representación $\lambda$ \\
$Y_\ell^m$ & Armónicos esféricos de grado $\ell$ y orden $m$ \\
$C_{l_1 m_1, l_2 m_2}^{l m}$ & Coeficientes de Clebsch-Gordan \\
$D^{(\ell)}(g)$ & Matriz de Wigner-D (representación irreducible de $SO(3)$ de grado $\ell$) \\
\end{longtable}

\subsection*{Álgebras de Clifford}
\begin{longtable}{>{\centering\arraybackslash}p{4cm}p{10cm}}
% Primera página
\toprule
\textbf{Símbolo} & \textbf{Descripción} \\
\midrule
\endfirsthead
% Páginas siguientes
\toprule
\textbf{Símbolo} & \textbf{Descripción} {\small(continuación)} \\
\midrule
\endhead
% Pie en páginas intermedias
\midrule
\multicolumn{2}{r}{\small\textit{Continúa en la siguiente página}} \\
\endfoot
% Pie en la última página
\bottomrule
\endlastfoot
$\text{Cl}(p,q)$ & Álgebra de Clifford con signatura $(p,q)$ \\
$\text{Cl}(3,0,1)$ & Álgebra geométrica proyectiva \\
$\text{Cl}(1,3)$ & Álgebra de Minkowski \\
$\mathbf{e}_i$ & Vectores base del álgebra \\
$a \wedge b$ & Producto exterior \\
$a \cdot b$ & Producto interior \\
$ab$ & Producto geométrico ($ab = a \cdot b + a \wedge b$) \\
$\tilde{a}$ & Reverso de $a$ \\
$|a|$ & Norma en el álgebra de Clifford \\
$\langle a \rangle_k$ & Parte de grado $k$ del multivector $a$ \\
$\eta_{ij}$ & Métrica del álgebra de Clifford \\
\end{longtable}

\section*{Geometría y Topología}

\subsection*{Variedades y Geometría Diferencial}
\begin{longtable}{>{\centering\arraybackslash}p{4cm}p{10cm}}
% Primera página
\toprule
\textbf{Símbolo} & \textbf{Descripción} \\
\midrule
\endfirsthead
% Páginas siguientes
\toprule
\textbf{Símbolo} & \textbf{Descripción} {\small(continuación)} \\
\midrule
\endhead
% Pie en páginas intermedias
\midrule
\multicolumn{2}{r}{\small\textit{Continúa en la siguiente página}} \\
\endfoot
% Pie en la última página
\bottomrule
\endlastfoot
$\mathcal{M}$ & Variedad diferenciable \\
$T_p\mathcal{M}$ & Espacio tangente a la variedad $\mathcal{M}$ en el punto $p$ \\
$T_p^*\mathcal{M}$ & Espacio cotangente (dual del espacio tangente) en el punto $p$ \\
$dF_p$ & Diferencial (pushforward) de la aplicación $F$ en el punto $p$ \\
$(df)_p$ & Diferencial de una función $f$ en $p$ (covector en $T_p^*\mathcal{M}$) \\
$dx^i$ & Elemento de la base dual del espacio cotangente (diferencial de la coordenada $x^i$) \\
$T\mathcal{M}$ & Fibrado tangente (unión disjunta de todos los espacios tangentes) \\
$\Gamma(TM)$ & Espacio de campos vectoriales en $\mathcal{M}$ \\
$\text{FM}(\mathcal{M})$ & Fibrado de marcos de la variedad $\mathcal{M}$ \\
$(U, \phi)$ & Carta local: $U \subset \mathcal{M}$ abierto, $\phi: U \to \mathbb{R}^m$ homeomorfismo \\
$g_{ij}$ & Componentes del tensor métrico en coordenadas locales \\
$g_p(v,w)$ & Producto interno de $v, w \in T_p\mathcal{M}$ inducido por la métrica \\
$\nabla_X Y$ & Derivada covariante del campo $Y$ en dirección $X$ \\
$\Gamma^k_{ij}$ & Símbolos de Christoffel de la conexión de Levi-Civita \\
$T(X,Y)$ & Tensor de torsión de una conexión: $T(X,Y) = \nabla_X Y - \nabla_Y X - [X,Y]$ \\
$R(X,Y)Z$ & Tensor de curvatura de Riemann \\
$\mathrm{Ric}$, $R_{ij}$ & Tensor de Ricci (contracción del tensor de Riemann) \\
$\mathcal{R}$ & Curvatura escalar ($\mathcal{R} = g^{ij}R_{ij}$) \\
$\exp_p: T_p\mathcal{M} \to \mathcal{M}$ & Mapa exponencial Riemanniano en el punto $p$ \\
$\log_p: \mathcal{M} \to T_p\mathcal{M}$ & Mapa logarítmico Riemanniano (inverso local de $\exp_p$) \\
$d_g(p,q)$ & Distancia geodésica entre puntos $p$ y $q$ \\
$\gamma: I \to \mathcal{M}$ & Curva en la variedad \\
$R \in SO(3)$ & Matriz de rotación en 3D \\
$\psi: S^2 \rightarrow \mathbb{R}^{d \times d'}$ & Kernel esférico para convolución en $S^2$ \\
$(P, \mathcal{M}, \pi, G)$ & Fibrado principal con espacio total $P$, base $\mathcal{M}$, proyección $\pi$ y grupo estructural $G$ \\
$\pi^{-1}(p)$ & Fibra sobre el punto $p \in \mathcal{M}$ \\
$A^*_u$ & Campo vectorial fundamental generado por $A \in \mathfrak{g}$ en el punto $u \in P$ \\
$\sigma_u: \mathfrak{g} \to V_u P$ & Isomorfismo canónico del álgebra de Lie al espacio vertical \\
$\omega \in \Omega^1(P, \mathfrak{g})$ & 1-forma de conexión en un fibrado principal (con valores en el álgebra de Lie) \\
$\Omega_\omega$ & 2-forma de curvatura de la conexión $\omega$ \\
$\Gamma_\gamma$ & Transporte paralelo a lo largo del camino $\gamma$ \\
$\text{Hol}_p(\omega)$ & Grupo de holonomía de la conexión $\omega$ en el punto $p$ \\
$K$ & Curvatura seccional (gaussiana en superficies) \\
\end{longtable}

\subsection*{Espacios Simétricos y Curvados}
\begin{longtable}{>{\centering\arraybackslash}p{4cm}p{10cm}}
% Primera página
\toprule
\textbf{Símbolo} & \textbf{Descripción} \\
\midrule
\endfirsthead
% Páginas siguientes
\toprule
\textbf{Símbolo} & \textbf{Descripción} {\small(continuación)} \\
\midrule
\endhead
% Pie en páginas intermedias
\midrule
\multicolumn{2}{r}{\small\textit{Continúa en la siguiente página}} \\
\endfoot
% Pie en la última página
\bottomrule
\endlastfoot
$\mathbb{R}^n$ & Espacio euclidiano $n$-dimensional \\
$\mathbb{Z}^d$ & Retículo entero $d$-dimensional \\
$S^n$ & Esfera $n$-dimensional \\
$\mathbb{RP}^{n-1}$ & Espacio proyectivo real $(n-1)$-dimensional \\
$\mathbb{CP}^{n-1}$ & Espacio proyectivo complejo $(n-1)$-dimensional \\
$T^n$ & Toro $n$-dimensional \\
$\text{Sym}^+(n)$ & Variedad de matrices simétricas positivas definidas $n \times n$ \\
$\mathbb{H}^n$ & Espacio hiperbólico $n$-dimensional \\
$\mathbb{D}^n$ & Disco de Poincaré (modelo conforme del espacio hiperbólico) \\
$d_{\mathbb{H}}(\mathbf{x}, \mathbf{y})$ & Distancia hiperbólica entre puntos \\
$\exp_{\mathbf{0}}^{\mathbb{H}}(\cdot)$ & Mapa exponencial hiperbólico en el origen \\
$\log_{\mathbf{0}}^{\mathbb{H}}(\cdot)$ & Mapa logarítmico hiperbólico en el origen \\
$\text{arcosh}(\cdot)$ & Función arcocoseno hiperbólico \\
$\langle \cdot, \cdot \rangle_{\mathcal{L}}$ & Producto interior de Lorentz \\
$\oplus_{\mathbb{H}}$ & Suma de Möbius en espacio hiperbólico \\
$\odot_{\mathbb{H}}$ & Multiplicación escalar de Möbius / agregación hiperbólica \\
\end{longtable}

\subsection*{Topología Algebraica}
\begin{longtable}{>{\centering\arraybackslash}p{4cm}p{10cm}}
% Primera página
\toprule
\textbf{Símbolo} & \textbf{Descripción} \\
\midrule
\endfirsthead
% Páginas siguientes
\toprule
\textbf{Símbolo} & \textbf{Descripción} {\small(continuación)} \\
\midrule
\endhead
% Pie en páginas intermedias
\midrule
\multicolumn{2}{r}{\small\textit{Continúa en la siguiente página}} \\
\endfoot
% Pie en la última página
\bottomrule
\endlastfoot
$\partial_k$ & Operador frontera de $k$-cadenas \\
$\partial_k^T$ & Operador cofrontera (adjunto del operador frontera) \\
$H_k(X)$ & $k$-ésimo grupo de homología \\
$H^k(X)$ & $k$-ésimo grupo de cohomología \\
$L_k$ & $k$-ésimo Laplaciano de Hodge \\
$\Omega^k$ & Espacio de $k$-formas diferenciales \\
$\sigma$ & Símplice \\
$\mathcal{K}$ & Complejo simplicial \\
$\text{pers}(b,d)$ & Persistencia de un par birth-death \\
\end{longtable}

\section*{Grafos y Estructuras Discretas}

\subsection*{Teoría de Grafos}
\begin{longtable}{>{\centering\arraybackslash}p{4cm}p{10cm}}
% Primera página
\toprule
\textbf{Símbolo} & \textbf{Descripción} \\
\midrule
\endfirsthead
% Páginas siguientes
\toprule
\textbf{Símbolo} & \textbf{Descripción} {\small(continuación)} \\
\midrule
\endhead
% Pie en páginas intermedias
\midrule
\multicolumn{2}{r}{\small\textit{Continúa en la siguiente página}} \\
\endfoot
% Pie en la última página
\bottomrule
\endlastfoot
$V, E$ & Vértices y aristas de un grafo \\
$\mathcal{G} = (V, E)$ & Grafo con vértices $V$ y aristas $E$ \\
$|V|$ & Número de vértices en el grafo \\
$|E|$ & Número de aristas en el grafo \\
$\mathcal{N}(v)$ & Vecindario del nodo $v$ en un grafo \\
$N_k(i)$ & Vecindario de tamaño $k$ del nodo $i$ (en grillas regulares) \\
$K_n$ & Grafo completo con $n$ nodos \\
$\vec{E}$ & Conjunto de aristas dirigidas \\
$(u, v) \in \vec{E}$ & Arista dirigida del vértice $u$ al vértice $v$ \\
$\text{DAG}$ & Grafo acíclico dirigido (Directed Acyclic Graph) \\
$\deg(v)$ & Grado del vértice $v$ \\
$A$ & Matriz de adyacencia \\
$D$ & Matriz de grados \\
$L$ & Matriz laplaciana del grafo ($L = D - A$) \\
$\tilde{L}$ & Matriz laplaciana normalizada \\
$\hat{f}$ & Transformada de Fourier en grafos de la señal $f$ \\
\end{longtable}

\section*{Redes Neuronales}

\subsection*{Arquitectura de Redes Neuronales}
\begin{longtable}{>{\centering\arraybackslash}p{4cm}p{10cm}}
% Primera página
\toprule
\textbf{Símbolo} & \textbf{Descripción} \\
\midrule
\endfirsthead
% Páginas siguientes
\toprule
\textbf{Símbolo} & \textbf{Descripción} {\small(continuación)} \\
\midrule
\endhead
% Pie en páginas intermedias
\midrule
\multicolumn{2}{r}{\small\textit{Continúa en la siguiente página}} \\
\endfoot
% Pie en la última página
\bottomrule
\endlastfoot
$\mathbf{W}$ & Matriz de pesos \\
$b$ & Vector de sesgos \\
$\boldsymbol{\theta}$ & Vector de parámetros del modelo \\
$\eta$ & Tasa de aprendizaje \\
$\sigma$ & Función de activación \\
$\tanh(\cdot)$ & Función tangente hiperbólica \\
$\text{ReLU}(\cdot)$ & Función de activación ReLU \\
$\text{softmax}(\cdot)$ & Función softmax \\
$\text{sigmoid}(\cdot)$ & Función sigmoide \\
$\mathcal{L}$ & Función de pérdida (loss) \\
$\mathcal{L}(\theta)$ & Función de pérdida parametrizada por $\theta$ \\
$\frac{\partial \mathcal{L}}{\partial \theta}$ & Gradiente de la pérdida respecto a parámetros \\
$\mathcal{N} = f^{(\ell)} \circ f^{(\ell-1)} \circ \cdots \circ f^{(1)}$ & Red neuronal como composición de capas \\
$\mathcal{G}_{\mathcal{N}} = (V_L, E_L)$ & Grafo de arquitectura de red (capas como vértices) \\
$\epsilon, \varepsilon$ & Parámetros de regularización \\
$\odot$ & Producto elemento a elemento (producto de Hadamard o Schur) \\
$\oslash$ & División elemento a elemento (división de Hadamard) \\
\end{longtable}

\subsection*{Parámetros de Capas}
\begin{longtable}{>{\centering\arraybackslash}p{4cm}p{10cm}}
% Primera página
\toprule
\textbf{Símbolo} & \textbf{Descripción} \\
\midrule
\endfirsthead
% Páginas siguientes
\toprule
\textbf{Símbolo} & \textbf{Descripción} {\small(continuación)} \\
\midrule
\endhead
% Pie en páginas intermedias
\midrule
\multicolumn{2}{r}{\small\textit{Continúa en la siguiente página}} \\
\endfoot
% Pie en la última página
\bottomrule
\endlastfoot
$c_{\text{in}}, c_{\text{out}}$ & Número de canales de entrada y salida \\
$N$ & Tamaño del batch (número de muestras) \\
$C$ & Número de canales en tensores \\
$H, W$ & Altura y anchura de tensores espaciales \\
$\Omega'$ & Dominio transformado (cuando difiere de $\Omega$) \\
$\alpha$ & Parámetro de pendiente en LeakyReLU \\
$n_{\text{classes}}$ & Número de clases en clasificación \\
$\gamma, \beta$ & Parámetros de escala y desplazamiento en normalización \\
$\mu$ & Media o estadística de centralización \\
$\sigma^2$ & Varianza (con subíndices $\mathcal{B}$, $L$, $G$, $IN$ según tipo de normalización) \\
$\mathcal{B}$ & Mini-batch de muestras \\
$m$ & Tamaño del mini-batch \\
$G$ & Número de grupos en Group Normalization \\
$p_{\text{drop}}$ & Probabilidad de dropout \\
$\mathbf{r}$ & Vector de máscaras binarias aleatorias (dropout) \\
$V_{\text{vocab}}$ & Tamaño del vocabulario (embeddings) \\
$E_{\text{emb}}$ & Matriz de embeddings \\
$d_{\text{emb}}$ & Dimensión del espacio de embeddings \\
$\mathbf{h}_t$ & Estado oculto en el tiempo $t$ (RNNs) \\
$\mathbf{x}_t$ & Entrada en el tiempo $t$ (secuencias) \\
$W_{\text{pool}}$ & Ventana de pooling \\
\end{longtable}

\subsection*{Redes Neuronales Convolucionales}
\begin{longtable}{>{\centering\arraybackslash}p{4cm}p{10cm}}
% Primera página
\toprule
\textbf{Símbolo} & \textbf{Descripción} \\
\midrule
\endfirsthead
% Páginas siguientes
\toprule
\textbf{Símbolo} & \textbf{Descripción} {\small(continuación)} \\
\midrule
\endhead
% Pie en páginas intermedias
\midrule
\multicolumn{2}{r}{\small\textit{Continúa en la siguiente página}} \\
\endfoot
% Pie en la última página
\bottomrule
\endlastfoot
$*$ & Operación de convolución estándar \\
$\star$ & Convolución en grupos (group convolution) \\
$K$ & Núcleo o filtro convolucional \\
$s$ & Stride (paso) de la convolución \\
$p$ & Padding (relleno) \\
$\text{pool}(\cdot)$ & Operación de pooling \\
$K_x, K_y$ & Kernels de gradiente (Sobel) en direcciones $x$ e $y$ \\
$(I * K)(x,y)$ & Operación de convolución discreta 2D \\
$\psi: \mathbb{Z}^2 \to \mathbb{R}$ & Filtro convolucional discreto \\
$\hat{f}(\omega)$ & Transformada de Fourier de la señal $f$ \\
$\mathcal{F}\{f\}$ & Operador de transformada de Fourier \\
$\mathcal{F}^{-1}\{F\}$ & Transformada de Fourier inversa \\
$\operatorname{sinc}(x)$ & Función sinc: $\sin(\pi x)/(\pi x)$ \\
$T_k(\lambda)$ & Polinomio de Chebyshev de grado $k$ \\
$\mathcal{R}[f]$ & Transformada de Radón de la función $f$ \\
\end{longtable}

\subsection*{Mecanismos de Atención}
\begin{longtable}{>{\centering\arraybackslash}p{4cm}p{10cm}}
% Primera página
\toprule
\textbf{Símbolo} & \textbf{Descripción} \\
\midrule
\endfirsthead
% Páginas siguientes
\toprule
\textbf{Símbolo} & \textbf{Descripción} {\small(continuación)} \\
\midrule
\endhead
% Pie en páginas intermedias
\midrule
\multicolumn{2}{r}{\small\textit{Continúa en la siguiente página}} \\
\endfoot
% Pie en la última página
\bottomrule
\endlastfoot
$\mathbf{Q}$ & Matriz de consultas (queries) \\
$\mathbf{K}$ & Matriz de claves (keys) \\
$\mathbf{V}$ & Matriz de valores (values) \\
$\mathbf{W}_Q, \mathbf{W}_K, \mathbf{W}_V$ & Matrices de pesos para Q, K, V \\
$\mathbf{W}_O$ & Matriz de proyección de salida \\
$d_k$ & Dimensión de las claves \\
$d_v$ & Dimensión de los valores \\
$d_{\text{model}}$ & Dimensión del modelo Transformer \\
$h$ & Número de cabezas de atención \\
$e_{ij}$ & Score de atención sin normalizar \\
$\alpha_{ij}$ & Peso de atención entre elementos $i$ y $j$ \\
$\boldsymbol{\alpha}_i$ & Vector de pesos de atención para consulta $i$ \\
$\mathbf{G}$ & Matriz de Gram en atención ($\mathbf{Q}\mathbf{K}^T$) \\
$\mathbf{S}$ & Matriz de scores de atención (sin normalizar) \\
$\text{Attention}(Q,K,V)$ & Mecanismo de atención básico \\
$\text{MultiHead}(Q,K,V)$ & Atención multi-cabeza \\
$\text{sim}(\mathbf{q}_i, \mathbf{k}_j)$ & Función de similaridad entre consulta $i$ y clave $j$ \\
$\mathcal{T}$ & Operador de atención (continuo o discreto según contexto) \\
$K(x,y)$ & Núcleo/kernel de atención entre puntos $x$ e $y$ \\
$\text{Spec}(A)$ & Espectro de la matriz $A$ (conjunto de eigenvalores) \\
$\operatorname{tr}(A)$ & Traza de la matriz $A$ \\
\end{longtable}

\section*{Geometric Deep Learning}

\subsection*{Marco Teórico del GDL}
\begin{longtable}{>{\centering\arraybackslash}p{4cm}p{10cm}}
% Primera página
\toprule
\textbf{Símbolo} & \textbf{Descripción} \\
\midrule
\endfirsthead
% Páginas siguientes
\toprule
\textbf{Símbolo} & \textbf{Descripción} {\small(continuación)} \\
\midrule
\endhead
% Pie en páginas intermedias
\midrule
\multicolumn{2}{r}{\small\textit{Continúa en la siguiente página}} \\
\endfoot
% Pie en la última página
\bottomrule
\endlastfoot
$\Omega$ & Dominio de datos (conjunto de posiciones donde se observan datos) \\
$\mathcal{S}$ & Estructura del dominio (relaciones topológicas, métricas o algebraicas) \\
$\tau$ & Topología (colección de conjuntos abiertos) \\
$\mathcal{D}$ & Geometría del dominio (par $(\Omega, \mathcal{S})$ de dominio y estructura) \\
$f$ & Señal (atributos observados sobre el dominio $\Omega$) \\
$f: \Omega \to \mathcal{C}$ & Señal como función del dominio al espacio de canales \\
$\mathcal{C}$ & Espacio de canales o atributos de la señal \\
$\mathcal{F}(\Omega, \mathcal{C})$ & Espacio de señales sobre $\Omega$ con valores en $\mathcal{C}$ \\
$(\Omega, \mathcal{S}, f)$ & Triple dominio-estructura-señal (representación completa de datos) \\
$(\mathcal{D}, f)$ & Par geometría-atributos (representación equivalente) \\
$\mathcal{X}, \mathcal{Y}$ & Espacios de entrada y salida \\
$\mathcal{F}$ & Espacio de características \\
$\mathcal{F}_\ell$ & Espacio de características en la capa $\ell$ \\
$\rho$ & Representación del grupo \\
$\phi$ & Función de mapeo de características \\
$\phi: X \rightarrow Y$ & Función equivariante entre espacios \\
$\psi$ & Función de agregación \\
$\oplus$ & Operación de agregación \\
\end{longtable}

\subsection*{Arquitecturas Equivariantes}
\begin{longtable}{>{\centering\arraybackslash}p{4cm}p{10cm}}
% Primera página
\toprule
\textbf{Símbolo} & \textbf{Descripción} \\
\midrule
\endfirsthead
% Páginas siguientes
\toprule
\textbf{Símbolo} & \textbf{Descripción} {\small(continuación)} \\
\midrule
\endhead
% Pie en páginas intermedias
\midrule
\multicolumn{2}{r}{\small\textit{Continúa en la siguiente página}} \\
\endfoot
% Pie en la última página
\bottomrule
\endlastfoot
$M_\ell$ & Función de mensaje en la capa $\ell$ \\
$U_\ell$ & Función de actualización en la capa $\ell$ \\
$m_v^{(\ell)}$ & Mensaje agregado al nodo $v$ en la capa $\ell$ \\
$\alpha_{ij}^{(\ell)}$ & Peso de atención entre nodos $i$ y $j$ en la capa $\ell$ \\
$Q^{(\ell)}, K^{(\ell)}, V^{(\ell)}$ & Matrices query, key, value en la capa $\ell$ \\
\end{longtable}
